\chapter{Trait \code{StringInternerBackend}}
\label{ch:interner-trait}

\begin{note}
Stub. Ver Capítulo~\ref{ch:interner} para el struct y las
métricas subyacentes. El trait existe como seam para un futuro
\code{GenerationalInterner} o \code{PartitionedInterner}; no se
entrega una segunda impl hoy, el trait está pre-posicionado
para cuando la evidencia de producción lo justifique.
\end{note}

\section{Contrato}

\begin{lstlisting}[style=rust]
pub trait StringInternerBackend {
    fn intern(&mut self, s: &str) -> StrId;
    fn resolve(&self, id: StrId) -> &str;
    fn try_resolve(&self, id: StrId) -> Option<&str>;
    fn get(&self, s: &str) -> Option<StrId>;
    fn len(&self) -> usize;
    fn is_empty(&self) -> bool;
    fn reserve(&mut self, additional: usize);
    fn metrics(&self) -> InternerMetrics;
}

pub struct InternerMetrics {
    pub unique_strings: usize,
    pub approx_bytes:   usize,
}
\end{lstlisting}

\section{Condiciones disparadoras}

Implementar un segundo backend vale la pena cuando cualquiera
de:
\begin{itemize}
    \item \code{metrics().unique\_strings} excede $\sim 10^7$ en
        una sesión viva.
    \item Una sesión streaming corre 24$\times$7 por más de 72h
        sin compactación.
    \item Un único proceso sirve a múltiples tenants
        (despliegue MSSP).
\end{itemize}

El motor expone las métricas vía
\code{api.get\_interner\_metrics()} para que un operador pueda
observar la aproximación a un gatillo en vez de chocarse con
él a ciegas.

\begin{inthecode}
\filepath{graph\_hunter\_core/src/interner.rs} --- trait +
métricas.\\
\filepath{TODO.md} --- criterios de disparo.
\end{inthecode}
